\documentclass[10pt,a4paper]{amsart}
\usepackage{amsmath,amssymb,amsthm,graphicx}
% Retained minimal preamble for the Stacks Project chapter varieties.tex.
% Verbatim excerpt of the chapter preamble: Xy-pic setup (native preamble lines 14-18), the native theorem
% declarations with the mathematical macros (native preamble lines 154-217), and the project own no-op
% definitions of the reference/slogan/history display environments (native wrapper.tex lines 1-3) that the
% counted statement uses for its slogan. Copied mechanically; nothing added, removed or redefined.
% For commutative diagrams we use Xy-pic
\usepackage[all]{xy}

% We use 2cell for 2-commutative diagrams.
\xyoption{2cell}
% Theorem environments.
%
\theoremstyle{plain}
\newtheorem{theorem}[subsection]{Theorem}
\newtheorem{proposition}[subsection]{Proposition}
\newtheorem{lemma}[subsection]{Lemma}

\theoremstyle{definition}
\newtheorem{definition}[subsection]{Definition}
\newtheorem{example}[subsection]{Example}
\newtheorem{exercise}[subsection]{Exercise}
\newtheorem{situation}[subsection]{Situation}

\theoremstyle{remark}
\newtheorem{remark}[subsection]{Remark}
\newtheorem{remarks}[subsection]{Remarks}

\numberwithin{equation}{subsection}

% Macros
%
\def\lim{\mathop{\mathrm{lim}}\nolimits}
\def\colim{\mathop{\mathrm{colim}}\nolimits}
\def\Spec{\mathop{\mathrm{Spec}}}
\def\Hom{\mathop{\mathrm{Hom}}\nolimits}
\def\Ext{\mathop{\mathrm{Ext}}\nolimits}
\def\SheafHom{\mathop{\mathcal{H}\!\mathit{om}}\nolimits}
\def\SheafExt{\mathop{\mathcal{E}\!\mathit{xt}}\nolimits}
\def\Sch{\mathit{Sch}}
\def\Mor{\mathop{\mathrm{Mor}}\nolimits}
\def\Ob{\mathop{\mathrm{Ob}}\nolimits}
\def\Sh{\mathop{\mathit{Sh}}\nolimits}
\def\NL{\mathop{N\!L}\nolimits}
\def\CH{\mathop{\mathrm{CH}}\nolimits}
\def\proetale{{pro\text{-}\acute{e}tale}}
\def\etale{{\acute{e}tale}}
\def\QCoh{\mathit{QCoh}}
\def\Ker{\mathop{\mathrm{Ker}}}
\def\Im{\mathop{\mathrm{Im}}}
\def\Coker{\mathop{\mathrm{Coker}}}
\def\Coim{\mathop{\mathrm{Coim}}}

% Boxtimes
%
\DeclareMathSymbol{\boxtimes}{\mathbin}{AMSa}{"02}

%
% Macros for moduli stacks/spaces
%
\def\QCohstack{\mathcal{QC}\!\mathit{oh}}
\def\Cohstack{\mathcal{C}\!\mathit{oh}}
\def\Spacesstack{\mathcal{S}\!\mathit{paces}}
\def\Quotfunctor{\mathrm{Quot}}
\def\Hilbfunctor{\mathrm{Hilb}}
\def\Curvesstack{\mathcal{C}\!\mathit{urves}}
\def\Polarizedstack{\mathcal{P}\!\mathit{olarized}}
\def\Complexesstack{\mathcal{C}\!\mathit{omplexes}}
% \Pic is the operator that assigns to X its picard group, usage \Pic(X)
% \Picardstack_{X/B} denotes the Picard stack of X over B
% \Picardfunctor_{X/B} denotes the Picard functor of X over B
\def\Pic{\mathop{\mathrm{Pic}}\nolimits}
\def\Picardstack{\mathcal{P}\!\mathit{ic}}
\def\Picardfunctor{\mathrm{Pic}}
\def\Deformationcategory{\mathcal{D}\!\mathit{ef}}
\newenvironment{reference}{}{}
\newenvironment{slogan}{}{}
\newenvironment{history}{}{}
\usepackage{cleveref}
\crefname{originalreference}{original reference}{original references}
\Crefname{originalreference}{Original reference}{Original references}
\makeatletter
\newlabel{coherent-lemma-flat-base-change-cohomology}{{02KH}{source}{}{sourcepointer-1}{}}
\newlabel{coherent-lemma-flat-base-change-cohomology@cref}{{[originalreference][0][]02KH}{[1][0][]source}{}{}{}}
\makeatother
\crefname{originalreference}{original reference}{original references}
\Crefname{originalreference}{Original reference}{Original references}
\makeatletter
\newlabel{algebra-proposition-ffdescent-finite-projectivity}{{058S}{source}{}{sourcepointer-2}{}}
\newlabel{algebra-proposition-ffdescent-finite-projectivity@cref}{{[originalreference][0][]058S}{[1][0][]source}{}{}{}}
\makeatother
\title{Faithful flat descent for a line bundle whose pullback is trivial: the sections form invertible modules and multiply isomorphically}
\author{The Stacks Project Authors}
\date{stacks-project:varieties.tex:local-snapshot-20261011; GNU FDL 1.2 or later}
\begin{document}
\maketitle
\noindent\textit{Editorial note.} All delivered mathematical text is the author's own, copied verbatim from the pinned Stacks Project chapter file varieties.tex; no statement or proof step was rewritten, reordered or removed.\par\medskip
\noindent\textit{Editorial note.} The counted span cites two cross-chapter labels, each recorded in assembly.external\_references with its canonical tag as a navigation pointer.\par\medskip
\noindent\textit{Editorial note.} The proof cites no bibliography item inside the counted span.\par\medskip
\noindent\textit{Allowed standard prerequisites (editorial).}\begin{itemize}
\item Faithful flat descent for finite projective (invertible) modules. Stacks Project, Algebra, Proposition algebra-proposition-ffdescent-finite-projectivity (tag 05B5). Concludes that M is invertible over R from invertibility after base change.
\item Flat base change for cohomology. Stacks Project, Cohomology of Schemes, Lemma coherent-lemma-flat-base-change-cohomology (tag 02KH). Identifies the sections of the base-changed sheaves.
\end{itemize}\medskip
\ifdefined\hypertarget\hypertarget{sourcepointer-1}{}\fi
\noindent\textit{Source reference.} Original reference 02KH: sources/stacks-chow/tags.txt, source lines 3372--3372. Cross-chapter standard prerequisite cited by the author inside the counted span; it is used as a pointer and its own proof is not part of this item.\par\medskip
\ifdefined\hypertarget\hypertarget{sourcepointer-2}{}\fi
\noindent\textit{Source reference.} Original reference 058S: sources/stacks-chow/tags.txt, source lines 6790--6790. Cross-chapter standard prerequisite cited by the author inside the counted span; it is used as a pointer and its own proof is not part of this item.\par\medskip
\begin{lemma}
\label{lemma-change-rings-pic-pre}
Let $A \to B$ be a faithfully flat ring map. Let $X$ be a quasi-compact and
quasi-separated scheme over $A$. Let $\mathcal{L}$ be an invertible
$\mathcal{O}_X$-module whose pullback to $X_B$ is trivial. Then
$H^0(X, \mathcal{L})$ and $H^0(X, \mathcal{L}^{\otimes -1})$ are invertible
$H^0(X, \mathcal{O}_X)$-modules and the
multiplication map induces an isomorphism
$$
H^0(X, \mathcal{L}) \otimes_{H^0(X, \mathcal{O}_X)}
H^0(X, \mathcal{L}^{\otimes -1}) \longrightarrow
H^0(X, \mathcal{O}_X)
$$
\end{lemma}

\begin{proof}
Denote $\mathcal{L}_B$ the pullback of $\mathcal{L}$ to $X_B$.
Choose an isomorphism $\mathcal{L}_B \to \mathcal{O}_{X_B}$.
Set $R = H^0(X, \mathcal{O}_X)$, $M = H^0(X, \mathcal{L})$ and think of
$M$ as an $R$-module. For every quasi-coherent $\mathcal{O}_X$-module
$\mathcal{F}$ with pullback $\mathcal{F}_B$ on $X_B$ there is a
canonical isomorphism
$H^0(X_B, \mathcal{F}_B) = H^0(X, \mathcal{F}) \otimes_A B$, see
Cohomology of Schemes, Lemma \ref{coherent-lemma-flat-base-change-cohomology}.
Thus we have
$$
M \otimes_R (R \otimes_A B) =
M \otimes_A B = H^0(X_B, \mathcal{L}_B) \cong
H^0(X_B, \mathcal{O}_{X_B}) = R \otimes_A B
$$
Since $R \to R \otimes_A B$ is faithfully flat (as the base change
of the faithfully flat map $A \to B$), we conclude
that $M$ is an invertible $R$-module by
Algebra, Proposition \ref{algebra-proposition-ffdescent-finite-projectivity}.
Similarly $N = H^0(X, \mathcal{L}^{\otimes -1})$ is an invertible $R$-module.
To see that the statement on tensor products is true, use that it is true
after pulling back to $X_B$ and faithful flatness of $R \to R \otimes_A B$.
Some details omitted.
\end{proof}

\end{document}
